\begin{definition}[Active median generator]
\label{def:philosophy-active-median-generator}
Let $M=\operatorname{Med}(\mathcal{R})$. A subset
$G\subseteq\mathcal{R}$ is an active median generator when
$$
  \operatorname{Med}(G)=M.
$$
It is minimal when no smaller subset of $\mathcal{R}$ has the same
median closure.
\end{definition}

\begin{theorem}[Minimal active median generators]
\label{thm:philosophy-minimal-active-median-generators}
The minimum size of an active median generator is $7$. There are exactly
two generators of this size:
$$
\begin{gathered}
  \{\mathrm{AGG},\mathrm{AUA},\mathrm{CUA},\mathrm{CUC},
    \mathrm{CUU},\mathrm{UAA},\mathrm{UCA}\},\\
  \{\mathrm{AGG},\mathrm{AUA},\mathrm{CUC},\mathrm{CUG},
    \mathrm{CUU},\mathrm{UAA},\mathrm{UCA}\}.
\end{gathered}
$$
\end{theorem}
\begin{proof}
Enumerate all subsets $G\subseteq\mathcal{R}$ and compute
$\operatorname{Med}(G)$ by iterating coordinatewise majority. No subset
of size at most $6$ generates $M$, and the two displayed subsets are
exactly the size-$7$ subsets that generate $M$.
\end{proof}

\begin{corollary}[Rigid median skeleton]
\label{cor:philosophy-rigid-median-skeleton}
The common part of the two minimum generators is
$$
  \{\mathrm{AGG},\mathrm{AUA},\mathrm{CUC},
    \mathrm{CUU},\mathrm{UAA},\mathrm{UCA}\}.
$$
The only choice is between $\mathrm{CUA}$ and $\mathrm{CUG}$.
\end{corollary}
\begin{proof}
Intersect the two generators in
\autoref{thm:philosophy-minimal-active-median-generators}.
\end{proof}

\begin{definition}[Median-essential codon]
\label{def:philosophy-median-essential-codon}
A codon $c\in\mathcal{R}$ is median-essential when
$$
  \operatorname{Med}(\mathcal{R}\setminus\{c\})
  \neq
  \operatorname{Med}(\mathcal{R}).
$$
Otherwise it is median-redundant.
\end{definition}

\begin{theorem}[Median-essential codons]
\label{thm:philosophy-median-essential-codons}
The median-essential codons are exactly
$$
  \mathrm{AGG},\quad \mathrm{AUA},\quad
  \mathrm{CUC},\quad \mathrm{CUU},\quad \mathrm{UCA}.
$$
The remaining active codons are median-redundant.
\end{theorem}
\begin{proof}
For each codon $c\in\mathcal{R}$ compute
$\operatorname{Med}(\mathcal{R}\setminus\{c\})$ and compare it with
$M$. The closure sizes are
$$
\begin{array}{c|c|c}
  c & |\operatorname{Med}(\mathcal{R}\setminus\{c\})| & \text{preserves }M\\
  \hline
  \mathrm{AAA} & 20 & \text{yes}\\
  \mathrm{AGA} & 20 & \text{yes}\\
  \mathrm{AGG} & 14 & \text{no}\\
  \mathrm{AUA} & 16 & \text{no}\\
  \mathrm{CUA} & 20 & \text{yes}\\
  \mathrm{CUC} & 19 & \text{no}\\
  \mathrm{CUG} & 20 & \text{yes}\\
  \mathrm{CUU} & 19 & \text{no}\\
  \mathrm{UAA} & 20 & \text{yes}\\
  \mathrm{UAG} & 20 & \text{yes}\\
  \mathrm{UCA} & 16 & \text{no}\\
  \mathrm{UGA} & 20 & \text{yes}\\
  \mathrm{UUA} & 20 & \text{yes}.
\end{array}
$$
\end{proof}

\begin{theorem}[Median deletion losses]
\label{thm:philosophy-median-deletion-losses}
The median-essential codons have the following closure losses:
$$
\begin{array}{c|l}
  c & M\setminus\operatorname{Med}(\mathcal{R}\setminus\{c\})\\
  \hline
  \mathrm{AGG} &
    \mathrm{AAG},\mathrm{ACG},\mathrm{AGG},
    \mathrm{AUG},\mathrm{UCG},\mathrm{UGG}\\
  \mathrm{AUA} &
    \mathrm{ACA},\mathrm{ACG},\mathrm{AUA},\mathrm{AUG}\\
  \mathrm{UCA} &
    \mathrm{ACA},\mathrm{ACG},\mathrm{UCA},\mathrm{UCG}\\
  \mathrm{CUC} &
    \mathrm{CUC}\\
  \mathrm{CUU} &
    \mathrm{CUU}.
\end{array}
$$
\end{theorem}
\begin{proof}
Subtract each deletion closure from $M$ for the five
median-essential codons of
\autoref{thm:philosophy-median-essential-codons}.
\end{proof}

\begin{corollary}[Weight and median generation differ]
\label{cor:philosophy-weight-median-generation-differ}
High reassignment weight is not the same as median-generating
importance. In particular $\mathrm{UGA}$, $\mathrm{UAA}$,
$\mathrm{UAG}$, $\mathrm{AGA}$, and $\mathrm{AAA}$ are
median-redundant, even though several of them have high reassignment
weight.
\end{corollary}
\begin{proof}
This is the redundant part of the deletion table in
\autoref{thm:philosophy-median-essential-codons}.
\end{proof}

\begin{theorem}[Complementary active pairs]
\label{thm:philosophy-complementary-active-pairs}
The reassigned support contains exactly two complementary pairs:
$$
  \{\mathrm{CUC},\mathrm{AGA}\},
  \qquad
  \{\mathrm{CUU},\mathrm{AGG}\}.
$$
Each pair has Hamming distance $6$ in $Q_{6}$.
\end{theorem}
\begin{proof}
Using $U=00$, $C=01$, $A=10$, and $G=11$, the bit words of
$\mathrm{CUC}$ and $\mathrm{AGA}$ are complementary, and the same holds
for $\mathrm{CUU}$ and $\mathrm{AGG}$. No other pair in
$\mathcal{R}$ has all six bits opposite.
\end{proof}

\begin{corollary}[Convex hull is too coarse]
\label{cor:philosophy-convex-hull-too-coarse}
The geodesic convex hull of $\mathcal{R}$ is all of $Q_{6}$, already
witnessed by either complementary pair in
\autoref{thm:philosophy-complementary-active-pairs}.
\end{corollary}
\begin{proof}
In a hypercube, the interval between complementary vertices is the
whole cube.
\end{proof}

\begin{definition}[$Q_{2}$ completion closure]
\label{def:philosophy-q2-completion-closure}
For $S\subseteq Q_{6}$, define $\operatorname{Comp}_{Q_{2}}(S)$ by
iterating the rule that whenever a $Q_{2}$ face contains at least three
current vertices, the whole face is adjoined.
\end{definition}

\begin{theorem}[Single-deletion completion robustness]
\label{thm:philosophy-single-deletion-completion-robustness}
For every active codon $c\in\mathcal{R}$,
$$
  \operatorname{Comp}_{Q_{2}}(\mathcal{R}\setminus\{c\})=Q_{6}.
$$
\end{theorem}
\begin{proof}
Delete each of the $13$ active codons and run the finite
$Q_{2}$-completion iteration. Every deletion case reaches all $64$
vertices.
\end{proof}

\begin{theorem}[Minimal percolating seeds]
\label{thm:philosophy-minimal-percolating-seeds}
Among subsets of $\mathcal{R}$, the minimum size of a
$Q_{2}$-completion seed that percolates to $Q_{6}$ is $7$. There are
$484$ such size-$7$ seeds.
\end{theorem}
\begin{proof}
Enumerate all subsets $G\subseteq\mathcal{R}$ and compute
$\operatorname{Comp}_{Q_{2}}(G)$. No subset of size at most $6$
percolates to $Q_{6}$, while $484$ subsets of size $7$ do.
\end{proof}

\begin{theorem}[Generator-closure robustness]
\label{thm:philosophy-generator-closure-robustness}
The reassigned support is rigid in median generation and robust in
$Q_{2}$-completion:
$$
\begin{gathered}
  \text{minimum active median generators: }2\text{ subsets of size }7,\\
  \text{minimum }Q_{2}\text{-percolating seeds: }484\text{ subsets of size }7.
\end{gathered}
$$
Thus the median closure identifies a small rigid generating skeleton,
whereas $Q_{2}$-completion records a broad percolation capacity.
\end{theorem}
\begin{proof}
Combine
\autoref{thm:philosophy-minimal-active-median-generators},
\autoref{thm:philosophy-single-deletion-completion-robustness}, and
\autoref{thm:philosophy-minimal-percolating-seeds}.
\end{proof}
